% ---------------------------------------------------------------------
% Argomento: diagramma di flusso con ciclo (niente vettori: un solo dato
%            letto, il ciclo aggiorna variabili con nome dato nel testo)
%   fila A: serbatoio riempito a secchi -> numero di secchi e livello finale
%   fila B: magazzino svuotato a spedizioni -> numero di spedizioni e pezzi rimasti
%   (niente tempo: il ciclo conta operazioni, non minuti)
%   fila C: scaffale riempito con una pila di scatole -> numero di scatole e altezza della pila
% Livello: intermedio (nessun suggerimento: condizione e aggiornamenti da ricavare)   |   Tempo stimato: 15 min
% ---------------------------------------------------------------------
\begin{esercizio}[12]
\ABC{%
Un serbatoio, inizialmente vuoto, si riempie versando un secchio d'acqua alla volta.
Disegna il diagramma di flusso di un algoritmo che conta quanti secchi servono per riempirlo.\\
\textbf{Dati da leggere:} $C$ = capacità del serbatoio (in litri); $S$ = litri contenuti in un secchio.\\
\textbf{Altre variabili:} $L$ = litri già presenti nel serbatoio; $N$ = numero di secchi versati.\\
Si versa un secchio dopo l'altro finché $L$ non raggiunge la capacità $C$. Alla fine visualizza $N$ e $L$.}{%
Un magazzino si svuota con una spedizione dopo l'altra. Disegna il diagramma di flusso di un
algoritmo che conta quante spedizioni si possono fare.\\
\textbf{Dati da leggere:} $Q$ = pezzi presenti in magazzino; $S$ = pezzi che partono con ogni spedizione.\\
\textbf{Altre variabili:} $K$ = numero di spedizioni fatte.\\
Si fanno spedizioni finché in magazzino ci sono abbastanza pezzi per una spedizione completa.
Alla fine visualizza $K$ e i pezzi rimasti $Q$.}{%
Su uno scaffale si impilano scatole tutte uguali, una sopra l'altra. Disegna il diagramma di
flusso di un algoritmo che conta quante scatole ci stanno.\\
\textbf{Dati da leggere:} $A$ = altezza dello scaffale (in cm); $H$ = altezza di una scatola (in cm).\\
\textbf{Altre variabili:} $P$ = altezza della pila; $N$ = numero di scatole impilate.\\
Si aggiungono scatole finché ce ne sta ancora una: la pila non deve mai superare l'altezza $A$.
Alla fine visualizza $N$ e $P$.}

\begin{solution}
\begin{center}
\ABC{%
\begin{tikzpicture}[x=1cm,y=1cm]
  \node[fcterm] (ini)  at (0,0) {Inizio};
  \node[fcio]   (in)   at (0,-1.1) {Leggi $C$, $S$};
  \node[fcinit] (ini0) at (0,-2.2) {$L \gets 0$, \ $N \gets 0$};
  \node[fcdec]  (d1)   at (0,-3.75) {$L < C$};
  \node[fcproc] (agg)  at (0,-5.3) {$L \gets L + S$};
  \node[fcproc] (inc)  at (0,-6.4) {$N \gets N + 1$};
  \node[fcio]   (out)  at (5.0,-3.75) {Scrivi $N$, $L$};
  \node[fcterm] (fin)  at (5.0,-4.95) {Fine};
  \draw[fcfreccia] (ini) -- (in);
  \draw[fcfreccia] (in) -- (ini0);
  \draw[fcfreccia] (ini0) -- (d1);
  \draw[fcfreccia] (d1) -- node[fcsi,right] {Sì} (agg);
  \draw[fcfreccia] (agg) -- (inc);
  \draw[fcfreccia] (inc.south) -- ++(0,-0.5) -| (-2.4,-2.95) -- (0,-2.95);
  \draw[fcfreccia] (d1) -- node[fcno,above] {No} (out);
  \draw[fcfreccia] (out) -- (fin);
\end{tikzpicture}}{%
\begin{tikzpicture}[x=1cm,y=1cm]
  \node[fcterm] (ini)  at (0,0) {Inizio};
  \node[fcio]   (in)   at (0,-1.1) {Leggi $Q$, $S$};
  \node[fcinit] (ini0) at (0,-2.2) {$K \gets 0$};
  \node[fcdec]  (d1)   at (0,-3.75) {$Q \geq S$};
  \node[fcproc] (agg)  at (0,-5.3) {$Q \gets Q - S$};
  \node[fcproc] (inc)  at (0,-6.4) {$K \gets K + 1$};
  \node[fcio]   (out)  at (5.0,-3.75) {Scrivi $K$, $Q$};
  \node[fcterm] (fin)  at (5.0,-4.95) {Fine};
  \draw[fcfreccia] (ini) -- (in);
  \draw[fcfreccia] (in) -- (ini0);
  \draw[fcfreccia] (ini0) -- (d1);
  \draw[fcfreccia] (d1) -- node[fcsi,right] {Sì} (agg);
  \draw[fcfreccia] (agg) -- (inc);
  \draw[fcfreccia] (inc.south) -- ++(0,-0.5) -| (-2.4,-2.95) -- (0,-2.95);
  \draw[fcfreccia] (d1) -- node[fcno,above] {No} (out);
  \draw[fcfreccia] (out) -- (fin);
\end{tikzpicture}}{%
\begin{tikzpicture}[x=1cm,y=1cm]
  \node[fcterm] (ini)  at (0,0) {Inizio};
  \node[fcio]   (in)   at (0,-1.1) {Leggi $A$, $H$};
  \node[fcinit] (ini0) at (0,-2.2) {$P \gets 0$, \ $N \gets 0$};
  \node[fcdec]  (d1)   at (0,-3.75) {$P + H \leq A$};
  \node[fcproc] (agg)  at (0,-5.3) {$P \gets P + H$};
  \node[fcproc] (inc)  at (0,-6.4) {$N \gets N + 1$};
  \node[fcio]   (out)  at (5.0,-3.75) {Scrivi $N$, $P$};
  \node[fcterm] (fin)  at (5.0,-4.95) {Fine};
  \draw[fcfreccia] (ini) -- (in);
  \draw[fcfreccia] (in) -- (ini0);
  \draw[fcfreccia] (ini0) -- (d1);
  \draw[fcfreccia] (d1) -- node[fcsi,right] {Sì} (agg);
  \draw[fcfreccia] (agg) -- (inc);
  \draw[fcfreccia] (inc.south) -- ++(0,-0.5) -| (-2.4,-2.95) -- (0,-2.95);
  \draw[fcfreccia] (d1) -- node[fcno,above] {No} (out);
  \draw[fcfreccia] (out) -- (fin);
\end{tikzpicture}}
\end{center}
\ABC{%
  Ogni giro del ciclo è un secchio versato: $N$ conta i giri. $L$ parte da 0 (serbatoio vuoto);
  quando il ciclo finisce, $L$ può superare di poco $C$ perché aumenta di $S$ alla volta.}{%
  Ogni giro del ciclo è una spedizione: $K$ conta i giri. Quando il ciclo finisce, in magazzino
  restano meno di $S$ pezzi (possono essere anche 0).}{%
  Ogni giro del ciclo è una scatola aggiunta: $N$ conta i giri. Prima di aggiungere una scatola si
  controlla che la pila con la scatola in più ($P+H$) non superi $A$; alla fine $P \leq A$.}
  \criterioABC
  {Inizio/Fine e simboli corretti (inizializzazione nell'esagono): 1 pt; lettura di $C$ e $S$: 1 pt;
    inizializzazione $L\gets 0$, $N\gets 0$: 2 pt; condizione del ciclo $L<C$ ricavata dal testo, con
    rami Sì/No giusti: 3 pt; aggiornamento $L\gets L+S$: 2 pt; contatore $N\gets N+1$: 1 pt; ritorno alla
    condizione: 1 pt; uscita di $N$ e $L$ dopo il ciclo: 1 pt. Ciclo con condizione in fondo
    (ripeti\ldots finché $L \geq C$) corretto: stessi punti.}
  {Inizio/Fine e simboli corretti (inizializzazione nell'esagono): 1 pt; lettura di $Q$ e $S$: 1 pt;
    inizializzazione $K\gets 0$ (e nessuna inizializzazione di $Q$ dopo la lettura): 2 pt; condizione del
    ciclo $Q\geq S$ ricavata da ``abbastanza pezzi per una spedizione completa'', con rami Sì/No giusti:
    3 pt ($Q>S$, che perde la spedizione quando $Q=S$: 1 pt); aggiornamento $Q\gets Q-S$: 2 pt; contatore $K\gets K+1$: 1 pt; ritorno
    alla condizione: 1 pt; uscita di $K$ e $Q$ dopo il ciclo: 1 pt. Ciclo con condizione in fondo
    corretto se gestisce il caso $Q<S$ iniziale: stessi punti.}
  {Inizio/Fine e simboli corretti (inizializzazione nell'esagono): 1 pt; lettura di $A$ e $H$: 1 pt;
    inizializzazione $P\gets 0$, $N\gets 0$: 2 pt; condizione del ciclo $P+H\leq A$ ricavata da ``ce ne
    sta ancora una'', con rami Sì/No giusti: 3 pt ($P<A$ o $P\leq A$, che fanno superare $A$: 1 pt);
    aggiornamento $P\gets P+H$: 2 pt; contatore $N\gets N+1$: 1 pt; ritorno alla condizione: 1 pt; uscita
    di $N$ e $P$ dopo il ciclo: 1 pt. Ciclo con condizione in fondo corretto se gestisce il caso $H>A$:
    stessi punti.}
\end{solution}
\end{esercizio}
